\section{Buffered low-energy spaces under norm perturbations}
\label{sec:buffered_spectral_norm}

A global operator-norm estimate gives both directed spectral comparisons,
provided that the target cutoff is enlarged. No restriction is imposed on
the spectrum between these cutoffs. The following auxiliary uses the
spectral-window argument of Arad, Landau, Vazirani and
Vidick~\cite[lines 955--968 and 1412--1420]{Arad2016RigorousRG}.
It retains the corrected finite window count recorded in
\cite{gap:alvv16_hard_truncation_windows}.

\begin{lemma}[Two directed buffered comparisons]
    \label{lem:buffered_spectral_norm_pair}
    \lean{Matrix.IsHermitian.isCloseTo_spectralSpanLE_pair_of_norm_sub_le}
    \leanok
    Let $H$ and $B$ be Hermitian matrices on a finite-dimensional complex
    Hilbert space. Suppose $\|H-B\|\leq\xi$ with $\xi>0$.
    For $c\in\mathbb R$ and $w>0$, put
    \begin{align}
        U&=\|H\|+\xi,&a&=\min\{c-1,-U\},\\
        n&=\left\lceil\frac{c-a}{2\xi}\right\rceil,&
        \delta&=\frac{n(2\xi)^2}{(w+\xi)^2}.
    \end{align}
    Write $E_A(c)$ for the spectral subspace of $A$ on $(-\infty,c]$.
    Then $E_B(c+w)$ is $\delta$-close to $E_H(c)$, and
    $E_H(c+w)$ is $\delta$-close to $E_B(c)$.
    Thus every vector $v\in E_H(c)$ satisfies
    \begin{align}
        \|P_{E_B(c+w)}v\|^2\geq(1-\delta)\|v\|^2,
    \end{align}
    and the corresponding assertion holds with $H$ and $B$ exchanged.
    The statement also allows the zero-dimensional Hilbert space.
\end{lemma}
\begin{proof}\leanok
    Both matrices have norm at most $U$, so all their eigenvalues are at
    least $-U\geq a$, while $a<c$. On every vector the perturbation has
    norm at most $\xi$ times the vector norm. Apply the finite-window
    comparison with window half-width $\xi$, lower endpoint $a$ and
    target cutoff $c+w$. Its numerical condition is exactly
    $n(2\xi)^2=\delta(w+\xi)^2$. Apply it a second time with the two
    matrices exchanged. The lower endpoint and loss are unchanged.
\end{proof}

The loss is selected from the original matrix and the error budget,
independently of the approximating matrix. If $0<\xi\leq1$, then
$n<\max\{1,c+U\}/(2\xi)+1$; in particular, fixed positive buffer and
polynomially bounded energy range require only inverse-polynomial norm
accuracy to make the loss inverse-polynomial. This auxiliary does not
assert a comparison at an unchanged cutoff or construct an efficient
filter.

\begin{lemma}[Stability of a separated low spectral cluster]
    \label{lem:spectral_cluster_norm_perturbation}
    \lean{Matrix.IsHermitian.spectral_cluster_of_norm_sub_le}
    \leanok
    Let $A,B$ be Hermitian matrices, let $a<b\le c<d$, and suppose
    $E_A(a)=E_A(d)$. Assume $\|A-B\|\le\xi$ with $\xi>0$.
    Put $U=\|A\|+\xi$ and
    \begin{align}
        n_-&=\left\lceil\frac{a-\min\{a-1,-U\}}{2\xi}\right\rceil,
        &\delta_-&=\frac{n_-(2\xi)^2}{(b-a+\xi)^2},\\
        n_+&=\left\lceil\frac{c-\min\{c-1,-U\}}{2\xi}\right\rceil,
        &\delta_+&=\frac{n_+(2\xi)^2}{(d-c+\xi)^2}.
    \end{align}
    If $\delta_-,\delta_+<1$, then
    \begin{align}
        E_B(b)&=E_B(c),&\dim E_B(b)&=\dim E_A(a).
    \end{align}
    The space $E_B(b)$ is $\delta_-$-close to $E_A(a)$, and
    $E_A(a)$ is $\delta_+$-close to $E_B(b)$. Every cutoff between
    $b$ and $c$ gives that same space and retains both directed estimates.
    Repeated eigenvalues, negative spectra and zero-dimensional spaces
    are allowed. This is an auxiliary for the approximating Hamiltonian
    in \cite[lines 1412--1420]{Arad2016RigorousRG}; the physical soft
    residuals and matrix approximation are separate constructions.
\end{lemma}
\begin{proof}\leanok
    \uses{lem:buffered_spectral_norm_pair,thm:subspace_close_nested_equal,
        thm:subspace_close_dimension}
    Apply the buffered comparison with source cutoff $a$ and buffer
    $b-a$, obtaining that $E_B(b)$ approximates $E_A(a)$ with loss
    $\delta_-$. Apply the reverse comparison with source cutoff $c$
    and buffer $d-c$, obtaining that $E_A(d)=E_A(a)$ approximates
    $E_B(c)$ with loss $\delta_+$. The two losses below one imply
    opposite dimension inequalities. Since $E_B(b)\subseteq E_B(c)$,
    the two spaces coincide and have the asserted dimension. Monotonicity
    gives equality and both estimates at every intermediate cutoff.
\end{proof}

\begin{lemma}[A uniform bound for the finite-window loss]
    \label{lem:spectral_window_uniform_envelope}
    \lean{Real.finite_window_loss_le_envelope,
        Real.finite_window_envelope_le,
        Real.finite_window_envelope_le_half_of_precision}
    \leanok
    Let $0\le A\le M$, $0\le c\le R$, and $\xi,w>0$. Put
    \begin{align}
        a&=\min\{c-1,-(A+\xi)\},&
        n&=\left\lceil\frac{c-a}{2\xi}\right\rceil,\notag\\
        n_*&=\left\lceil\frac{R+M+\xi+1}{2\xi}\right\rceil,&
        \Lambda&=\frac{n_*(2\xi)^2}{(w+\xi)^2}.
    \end{align}
    Then $n(2\xi)^2/(w+\xi)^2\le\Lambda$. If $\xi\le1$, then
    \begin{align}
        \Lambda\le\frac{2\xi(R+M+4)}{w^2}.
        \label{eq:spectral_window_linear_precision}
    \end{align}
    Consequently, $\xi\le\beta w^2/[4(R+M+4)]$ implies
    $\Lambda\le\beta/2$.
\end{lemma}
\begin{proof}\leanok
    The lower endpoint gives
    $c-a=\max\{1,c+A+\xi\}\le R+M+\xi+1$, so $n\le n_*$.
    The ceiling bound $n_*\le(R+M+\xi+1)/(2\xi)+1$ and
    $\xi^2\le\xi$ yield
    \begin{align}
        n_*(2\xi)^2
          &\le2\xi(R+M+1)+6\xi^2
           \le2\xi(R+M+4).
    \end{align}
    Dividing by $(w+\xi)^2\ge w^2>0$ proves
    (\ref{eq:spectral_window_linear_precision}). Substitute the final
    precision inequality to obtain $\beta/2$.
\end{proof}

\begin{corollary}[A positive common precision]
    \label{cor:spectral_window_precision_choice}
    \lean{Real.finite_window_precision_parameters,
        Real.soft_cluster_window_precision_parameters}
    \leanok
    Let $M,R\ge0$ and $w,\beta>0$. The choice
    \begin{align}
        \xi=\min\left\{1,\frac{\beta w^2}{4(R+M+4)}\right\}>0
    \end{align}
    satisfies both precision inequalities of
    Lemma~\ref{lem:spectral_window_uniform_envelope}. Thus its uniform
    loss is at most $\beta/2$, before $A$ and $c$ are specified.
    In particular, for $M\ge0$ and $g,\tau>0$, the values
    $R=4+g$, $w=g/5$, and $\beta=\tau/200$ give
    \begin{align}
        \xi=\min\left\{1,
          \frac{\tau g^2}{20000(M+g+8)}\right\},
        \qquad\Lambda\le\frac{\tau}{400}.
    \end{align}
    The same precision therefore controls every cutoff $0\le c\le4+g$
    and every norm bound $0\le A\le M$.
\end{corollary}
\begin{proof}\leanok
    \uses{lem:spectral_window_uniform_envelope}
    Both arguments of the minimum are positive. The minimum is at most
    one and at most the displayed quotient, so the lemma applies.
    The specialization follows from
    $(\tau/200)(g/5)^2/[4(4+g+M+4)]
      =\tau g^2/[20000(M+g+8)]$.
\end{proof}
